\documentclass{amsart}
\usepackage[T1]{fontenc}
\usepackage{amsmath,amssymb,amsthm}
\usepackage{mathtools}
\usepackage{natbib}

\newtheorem{theorem}{Theorem}[section]
\newtheorem{proposition}[theorem]{Proposition}
\newtheorem{lemma}[theorem]{Lemma}
\newtheorem{corollary}[theorem]{Corollary}
\theoremstyle{definition}
\newtheorem{definition}[theorem]{Definition}
\newtheorem{remark}[theorem]{Remark}

\newcommand{\norm}[1]{\left\lVert #1 \right\rVert}
\newcommand{\abs}[1]{\left\lvert #1 \right\rvert}
\newcommand{\ip}[2]{\langle #1, #2 \rangle}
\newcommand{\ext}{\operatorname{ext}}
\newcommand{\co}{\operatorname{co}}

\title{Extreme Points and the Krein--Milman Theorem\\
Notes on Brezis Problem 1}

\author{Nam Le}
\address{Department of Computer Science,
    Faculty of Information Technology, University of Science, Vietnam National University,
    227 Nguyen Van Cu Street, Cho Quan Ward, Ho Chi Minh City}
\email{lnnam@hcmus.edu.vn}
\date{\today}

\begin{document}

\begin{abstract}
We organize a complete solution to Problem~1 from Brezis' \emph{Functional Analysis,
Sobolev Spaces and Partial Differential Equations}. The argument is arranged as a short
expository note: we begin with elementary characterizations of extreme points, prove that
every nonempty extreme subset of a compact convex set contains an extreme point, derive the
Krein--Milman theorem, and conclude with explicit descriptions of the extreme points of the
closed unit ball in several classical Banach spaces.
\end{abstract}

\maketitle

\section{Introduction}
The Krein--Milman theorem is one of the structural results at the heart of convexity in
functional analysis: a compact convex set in a locally convex setting is generated, after
taking closed convex hulls, by its extreme points. Problem~1 in Brezis' text isolates the
main ideas behind this theorem and turns them into a sequence of manageable lemmas
\citep{brezis2011}.

This note follows the order of the exercise, but we reorganize it into a proof-driven
narrative. We first collect the basic equivalences for extreme points and extreme sets. We
then prove that every nonempty extreme subset of a compact convex set contains an extreme
point; this is the key step that allows one to deduce the Krein--Milman theorem by a
separation argument. Finally, we compute the extreme points of the closed unit ball in a
number of standard Banach spaces, thereby illustrating how geometry depends strongly on the
ambient norm.

\section{Preliminaries}
Let $E$ be a normed vector space and let $K \subset E$ be convex.

\begin{definition}
A point $a \in K$ is an \emph{extreme point} of $K$ if whenever
\[
a = tx + (1-t)y, \qquad t \in (0,1), \quad x,y \in K,
\]
one necessarily has $x = y = a$.
\end{definition}

When $K$ is assumed compact and convex, a nonempty closed subset $M \subset K$ is called an
\emph{extreme set} if for every $t \in (0,1)$,
\[
tx + (1-t)y \in M, \quad x,y \in K \implies x,y \in M.
\]
Thus extreme sets are the closed faces of $K$.

We shall also use the Hahn--Banach theorem in its elementary separation form: if $x \neq y$
in a normed vector space, then there exists $f \in E^*$ such that $\ip{f}{x} \neq \ip{f}{y}$.
This follows by defining a continuous linear functional on $\operatorname{span}\{x-y\}$ and
extending it to all of $E$.

\section{Extreme Points and Extreme Sets}

\begin{proposition}\label{prop:characterization}
Let $a \in K$. Then $a$ is an extreme point of $K$ if and only if $K \setminus \{a\}$ is
convex.
\end{proposition}

\begin{proof}
Assume first that $a$ is an extreme point, and let $x,y \in K \setminus \{a\}$ with
$t \in (0,1)$. Since $K$ is convex, $tx + (1-t)y \in K$. If this convex combination were equal
to $a$, then the extremality of $a$ would force $x = y = a$, which is impossible. Hence
$tx + (1-t)y \in K \setminus \{a\}$, so $K \setminus \{a\}$ is convex.

Conversely, assume that $K \setminus \{a\}$ is convex, and suppose
$a = tx + (1-t)y$ with $t \in (0,1)$ and $x,y \in K$. If $x \neq a$ and $y \neq a$, then
$x,y \in K \setminus \{a\}$, hence the convexity of $K \setminus \{a\}$ implies
$a = tx + (1-t)y \in K \setminus \{a\}$, a contradiction. Therefore at least one of $x,y$ is
equal to $a$, and then the identity $a = tx + (1-t)y$ with $t \in (0,1)$ forces the other one to
be equal to $a$ as well. Thus $a$ is extreme.
\end{proof}

\begin{proposition}\label{prop:finite-combination}
Let $a$ be an extreme point of $K$. Let $x_1,\dots,x_n \in K$ and let
$\alpha_1,\dots,\alpha_n > 0$ satisfy $\sum_{i=1}^n \alpha_i = 1$ and
$\sum_{i=1}^n \alpha_i x_i = a$. Then $x_i = a$ for every $i$.
\end{proposition}

\begin{proof}
We argue by induction on $n$. The case $n=2$ is exactly the definition of an extreme point.
Assume the statement true for $n-1$ terms and write
\[
a = \alpha_1 x_1 + (1-\alpha_1)z, \qquad
z = \sum_{i=2}^n \frac{\alpha_i}{1-\alpha_1} x_i \in K.
\]
Since $a$ is extreme and $\alpha_1 \in (0,1)$, we obtain $x_1 = a$ and $z = a$. The induction
hypothesis applied to the representation of $z$ then yields $x_i = a$ for $i=2,\dots,n$.
\end{proof}

\begin{proposition}\label{prop:singleton-face}
Let $a \in K$, and assume now that $K$ is compact and convex. Then $a$ is an extreme point of
$K$ if and only if $\{a\}$ is an extreme set of $K$.
\end{proposition}

\begin{proof}
If $a$ is an extreme point, then $\{a\}$ is nonempty and closed. If
$tx + (1-t)y \in \{a\}$ for some $t \in (0,1)$ and $x,y \in K$, then
$a = tx + (1-t)y$, hence $x = y = a$. Therefore $\{a\}$ is an extreme set.

Conversely, if $\{a\}$ is an extreme set and $a = tx + (1-t)y$ with $t \in (0,1)$ and
$x,y \in K$, then $tx + (1-t)y \in \{a\}$, so the definition of an extreme set gives
$x,y \in \{a\}$. Hence $x = y = a$, and $a$ is extreme.
\end{proof}

\section{Every Extreme Set Contains an Extreme Point}
From now on, $K \subset E$ is assumed to be nonempty, compact, and convex.

\begin{proposition}\label{prop:support-face}
Let $A \subset K$ be an extreme set and let $f \in E^*$. Define
\[
B = \left\{x \in A : \ip{f}{x} = \max_{y \in A} \ip{f}{y} \right\}.
\]
Then $B$ is an extreme set of $K$.
\end{proposition}

\begin{proof}
Because $A$ is compact and $f$ is continuous, the maximum is attained, so $B$ is nonempty.
Since $B = A \cap f^{-1}(\{\max_A f\})$, it is closed.

Let $t \in (0,1)$ and suppose $tx + (1-t)y \in B$ for some $x,y \in K$. Since $B \subset A$ and
$A$ is an extreme set, we first obtain $x,y \in A$. Let $m = \max_{z \in A} \ip{f}{z}$. Then
\[
m = \ip{f}{tx + (1-t)y} = t\ip{f}{x} + (1-t)\ip{f}{y} \le tm + (1-t)m = m.
\]
Hence equality must hold in both inequalities, so $\ip{f}{x} = \ip{f}{y} = m$, which means
$x,y \in B$. Thus $B$ is an extreme set.
\end{proof}

\begin{proposition}\label{prop:zorn}
Let $M \subset K$ be an extreme set. Consider the family $\mathcal{F}$ of all extreme sets of $K$
contained in $M$, ordered by
\[
A \leq B \quad \Longleftrightarrow \quad B \subset A.
\]
Then $\mathcal{F}$ has a maximal element.
\end{proposition}

\begin{proof}
The set $\mathcal{F}$ is nonempty because $M \in \mathcal{F}$. Let $\mathcal{C} \subset \mathcal{F}$ be a
chain. Because the order is reverse inclusion, the members of $\mathcal{C}$ are nested by
inclusion. Set
\[
A_\mathcal{C} = \bigcap_{A \in \mathcal{C}} A.
\]
Each $A \in \mathcal{C}$ is a nonempty closed subset of the compact set $M$, and the chain
condition implies the finite intersection property. Therefore $A_\mathcal{C}$ is nonempty and
closed. We claim that it is an extreme set. Indeed, if $tx + (1-t)y \in A_\mathcal{C}$ with
$t \in (0,1)$ and $x,y \in K$, then $tx + (1-t)y \in A$ for every $A \in \mathcal{C}$, so the
extremality of each $A$ yields $x,y \in A$ for every $A \in \mathcal{C}$. Thus
$x,y \in A_\mathcal{C}$.

Since $A_\mathcal{C} \subset A$ for all $A \in \mathcal{C}$, we have $A \leq A_\mathcal{C}$ for all
$A \in \mathcal{C}$. Thus every chain has an upper bound, and Zorn's lemma gives a maximal
element $M_0 \in \mathcal{F}$.
\end{proof}

\begin{proposition}\label{prop:singleton}
The maximal element $M_0$ furnished by Proposition~\ref{prop:zorn} is reduced to a single point.
\end{proposition}

\begin{proof}
Assume, by contradiction, that $M_0$ contains two distinct points $x$ and $y$. By the
Hahn--Banach theorem there exists $f \in E^*$ such that $\ip{f}{x} \neq \ip{f}{y}$, so $f$ is not
constant on $M_0$. Let
\[
B = \left\{z \in M_0 : \ip{f}{z} = \max_{w \in M_0} \ip{f}{w} \right\}.
\]
By Proposition~\ref{prop:support-face}, $B$ is an extreme set of $K$ contained in $M_0$.
Moreover $B$ is a proper subset of $M_0$, because $f$ is not constant on $M_0$. Therefore
$M_0 < B$ in the reverse-inclusion order, contradicting the maximality of $M_0$. Hence $M_0$
must be a singleton.
\end{proof}

\begin{corollary}\label{cor:extreme-set-contains-extreme-point}
Every nonempty extreme set of $K$ contains at least one extreme point of $K$.
\end{corollary}

\begin{proof}
Let $M \subset K$ be a nonempty extreme set. By Proposition~\ref{prop:zorn}, the family of
extreme subsets of $K$ contained in $M$ has a maximal element $M_0$, and by
Proposition~\ref{prop:singleton} we may write $M_0 = \{a\}$. Proposition~\ref{prop:singleton-face}
then shows that $a$ is an extreme point of $K$.
\end{proof}

\section{The Krein--Milman Theorem}

\begin{theorem}[Krein--Milman for compact convex subsets of normed spaces]
Let $K$ be a nonempty compact convex subset of $E$. Then
\[
K = \overline{\co}(\ext K).
\]
\end{theorem}

\begin{proof}
Set $C = \overline{\co}(\ext K)$. By definition, $C$ is a closed convex subset of $K$. Suppose that
$C \neq K$, and choose $x_0 \in K \setminus C$. Since $C$ is closed and convex, the Hahn--Banach
separation theorem provides $f \in E^*$ such that
\[
\sup_{x \in C} \ip{f}{x} < \ip{f}{x_0}.
\]

Let
\[
A = \left\{x \in K : \ip{f}{x} = \max_{y \in K} \ip{f}{y} \right\}.
\]
By Proposition~\ref{prop:support-face}, $A$ is an extreme set of $K$. Therefore,
Corollary~\ref{cor:extreme-set-contains-extreme-point} yields an extreme point $a \in A$. Since
$a \in \ext K \subset C$, we have
\[
\ip{f}{a} \le \sup_{x \in C} \ip{f}{x} < \ip{f}{x_0}.
\]
On the other hand, $a \in A$ means that $\ip{f}{a} = \max_{y \in K} \ip{f}{y} \ge \ip{f}{x_0}$,
a contradiction. Hence $C = K$.
\end{proof}

\begin{remark}
The theorem is usually stated in the larger framework of locally convex spaces. In the present
exercise, the normed-space setting already contains the essential idea: compactness gives
supporting faces, and Hahn--Banach supplies the separating functionals.
\end{remark}

\section{Extreme Points of Unit Balls in Classical Spaces}
We now determine the set of extreme points of the closed unit ball $B_E$ in the six spaces
listed in Brezis' problem.

\begin{proposition}\label{prop:linfty}
The extreme points of the closed unit ball of $\ell^\infty$ are exactly the sequences
$x = (x_n)$ such that $\abs{x_n} = 1$ for every $n$.
\end{proposition}

\begin{proof}
Assume first that $x \in B_{\ell^\infty}$ and that $\abs{x_k} < 1$ for some index $k$. Choose
$\varepsilon > 0$ such that $\abs{x_k} + \varepsilon \le 1$, and define
\[
y = x + \varepsilon e_k, \qquad z = x - \varepsilon e_k.
\]
Then $y \neq z$, $x = \tfrac12(y+z)$, and $\norm{y}_\infty, \norm{z}_\infty \le 1$, so $x$ is not
extreme.

Conversely, assume that $\abs{x_n} = 1$ for every $n$ and write $x = \tfrac12(y+z)$ with
$\norm{y}_\infty, \norm{z}_\infty \le 1$. For each $n$ we have
\[
2x_n = y_n + z_n.
\]
If $x_n = 1$, then $y_n + z_n = 2$ with $y_n,z_n \le 1$, hence $y_n = z_n = 1$. Similarly, if
$x_n = -1$, then $y_n = z_n = -1$. Thus $y = z = x$, and $x$ is extreme.
\end{proof}

\begin{proposition}\label{prop:c}
The extreme points of the closed unit ball of $c$ are exactly the convergent sign sequences,
that is,
\[
\ext(B_c) = \{x = (x_n) \in c : x_n \in \{-1,1\} \text{ for all } n\}.
\]
Equivalently, these are the $\{-1,1\}$-valued sequences that are eventually constant.
\end{proposition}

\begin{proof}
As a closed subspace of $\ell^\infty$, the space $c$ inherits the same perturbation argument as in
Proposition~\ref{prop:linfty}. Thus any extreme point must satisfy $\abs{x_n} = 1$ for all $n$.
Conversely, if $x \in c$ is such a sequence and $x = \tfrac12(y+z)$ with $y,z \in B_c$, the
coordinate-wise argument from Proposition~\ref{prop:linfty} again gives $y = z = x$. Hence $x$
is extreme. Since a convergent sequence taking only the values $\pm 1$ must eventually be
constant, the final description follows.
\end{proof}

\begin{proposition}\label{prop:c0}
The closed unit ball of $c_0$ has no extreme points.
\end{proposition}

\begin{proof}
Let $x \in B_{c_0}$. Because $x_n \to 0$, there exists some index $k$ such that $\abs{x_k} < 1$.
The perturbation argument used in Proposition~\ref{prop:linfty} then shows that $x$ is the
midpoint of two distinct points of $B_{c_0}$. Therefore $x$ is not extreme. Since $x$ was
arbitrary, $\ext(B_{c_0}) = \varnothing$.
\end{proof}

\begin{proposition}\label{prop:l1}
The extreme points of the closed unit ball of $\ell^1$ are exactly the vectors $\pm e_n$,
$n \ge 1$.
\end{proposition}

\begin{proof}
We first prove that each $\pm e_n$ is extreme. Suppose
$e_n = \tfrac12(y+z)$ with $\norm{y}_1, \norm{z}_1 \le 1$. Then
\[
1 = \norm{e_n}_1 \le \tfrac12(\norm{y}_1 + \norm{z}_1) \le 1,
\]
so equality holds throughout. In particular, $\norm{y}_1 = \norm{z}_1 = 1$. Since the $n$th
coordinate of $e_n$ is $1$, we have $y_n + z_n = 2$ with $\abs{y_n}, \abs{z_n} \le 1$, hence
$y_n = z_n = 1$. The equality $\norm{y}_1 = 1$ then forces all other coordinates of $y$ to be
zero, and similarly for $z$. Thus $y = z = e_n$. The argument for $-e_n$ is identical.

Now let $x \in B_{\ell^1}$ and assume that $x \neq \pm e_n$ for all $n$. If $\norm{x}_1 < 1$, choose
$n$ and $\varepsilon > 0$ with $\norm{x}_1 + \varepsilon \le 1$; then $x$ is the midpoint of
$x \pm \varepsilon e_n$, both belonging to $B_{\ell^1}$. Hence $x$ is not extreme.

It remains to consider the case $\norm{x}_1 = 1$. Since $x$ is not one of the vectors $\pm e_n$,
there exist distinct indices $i,j$ such that $x_i \neq 0$ and $x_j \neq 0$. Choose
$0 < \varepsilon < \min\{\abs{x_i}, \abs{x_j}\}$ and define $h \in \ell^1$ by
\[
h_i = \operatorname{sgn}(x_i), \qquad h_j = -\operatorname{sgn}(x_j), \qquad h_k = 0 \text{ for }
k \notin \{i,j\}.
\]
Then the signs of the $i$th and $j$th coordinates do not change in $x \pm \varepsilon h$, and a
direct computation gives
\[
\norm{x \pm \varepsilon h}_1 = \norm{x}_1 = 1.
\]
Since $x = \tfrac12((x+\varepsilon h) + (x-\varepsilon h))$ with $x+\varepsilon h \neq x-\varepsilon h$,
$x$ is not extreme. We conclude that $\ext(B_{\ell^1}) = \{\pm e_n : n \ge 1\}$.
\end{proof}

\begin{proposition}\label{prop:lp}
Let $1 < p < \infty$. The extreme points of the closed unit ball of $\ell^p$ are exactly the unit
sphere:
\[
\ext(B_{\ell^p}) = \{x \in \ell^p : \norm{x}_p = 1\}.
\]
\end{proposition}

\begin{proof}
If $\norm{x}_p < 1$, then $x$ is an interior point of $B_{\ell^p}$, hence certainly not extreme.

Conversely, assume $\norm{x}_p = 1$ and write $x = \tfrac12(y+z)$ with
$\norm{y}_p, \norm{z}_p \le 1$. Since the function $t \mapsto \abs{t}^p$ is strictly convex for
$1 < p < \infty$, the norm on $\ell^p$ is strictly convex. Therefore equality in
\[
1 = \norm{x}_p \le \tfrac12\norm{y}_p + \tfrac12\norm{z}_p \le 1
\]
forces $y = z = x$. Thus every point of the unit sphere is extreme.
\end{proof}

\begin{proposition}\label{prop:l1R}
The closed unit ball of $L^1(\mathbb{R})$ has no extreme points.
\end{proposition}

\begin{proof}
Let $f \in B_{L^1(\mathbb{R})}$. If $\norm{f}_1 < 1$, then $f$ is an interior point of the unit ball,
so it is not extreme.

Assume now that $\norm{f}_1 = 1$. Consider the finite measure
\[
\nu(A) = \int_A \abs{f(x)}\,dx.
\]
Because Lebesgue measure is nonatomic, so is $\nu$. Hence there exists a measurable set $A$
such that $\nu(A) = \tfrac12$. Define
\[
g = 2f\mathbf{1}_A, \qquad h = 2f\mathbf{1}_{\mathbb{R} \setminus A}.
\]
Then
\[
\norm{g}_1 = 2\nu(A) = 1, \qquad \norm{h}_1 = 2\nu(\mathbb{R} \setminus A) = 1,
\]
and
\[
f = \tfrac12(g+h).
\]
Since $\nu(A) = \nu(\mathbb{R} \setminus A) = \tfrac12$, both $g$ and $h$ are nonzero, hence
$g \neq h$. Therefore $f$ is not extreme. We conclude that $\ext(B_{L^1(\mathbb{R})}) = \varnothing$.
\end{proof}

\section{Conclusion}
Problem~1 packages the proof of the Krein--Milman theorem into a sequence of elementary but
revealing reductions. The central mechanism is the passage from a compact convex set to a
smaller support face on which a separating functional is maximized. Once one knows that every
nonempty extreme subset contains an extreme point, the final separation argument becomes
natural and transparent. The examples of unit balls show that this abstract theorem is not
merely existential: the geometry of extreme points can vary from a rich family of sign
sequences in $\ell^\infty$ to complete absence of extreme points in $c_0$ and $L^1(\mathbb{R})$.

\bibliographystyle{plainnat}
\bibliography{references}

\end{document}